% ====================================================================
\chapter{Number Systems and the Properties of Real Numbers}\label{ch:numbers}
% ====================================================================

\begin{chapterintro}
Economics measures prices, quantities, incomes and rates, and all of these are numbers. This chapter
introduces the four number systems used in the module --- the natural numbers, the integers, the rational
numbers and the real numbers --- and then lists the basic properties of the real numbers from which every
calculation in this book can be justified. These properties are the axioms of the module: when we solve an
inequality, every step is an application of one of them.
\end{chapterintro}

\begin{prerequisites}
Sets, subsets and set-builder notation (\cref{ch:sets}); the meaning of axiom and theorem
(\cref{sec:proof}); implication and equivalence (\cref{sec:implication,sec:equivalence}).
\end{prerequisites}

\section{The four number systems}\label{sec:number-systems}

\begin{definition}[Number systems][def:number-systems]
\begin{itemize}
  \item The \term{natural numbers} are the set $\N=\{1,2,3,\dots\}$.
  \item The \term{integers} are the set $\Z=\{\dots,-3,-2,-1,0,1,2,3,\dots\}$.
  \item The \term{rational numbers} are the set $\Q$ consisting of all numbers that can be written as $p/q$
  with $p\in\Z$ and $q\in\N$.
  \item The \term{real numbers} are (informally) the set $\R$ of points on the number line.
\end{itemize}
\src{L1, slide 16}
\end{definition}

\begin{core}[Successor property of $\N$]
For any $n\in\N$, we have $n+1\in\N$. \src{L1, slide 16}

Starting from $1$ and repeatedly adding $1$ reaches every natural number and never leaves $\N$. This simple
fact is what makes proof by induction work (\cref{ch:ind}).
\end{core}

\begin{warning}[$0$ is not a natural number in this module]
In ECN115, $\N=\{1,2,3,\dots\}$ starts at $1$. Some books include $0$. The difference matters for index
ranges in sums and for the starting point of an induction. When $0$ is wanted it is added explicitly, as in
$\{0\}\cup\N$ in the binomial formula (\cref{ch:binom}).
\end{warning}

\paragraph{Why $q\in\N$ in the definition of $\Q$} Requiring the denominator to be a natural number does
two jobs at once. It rules out $q=0$, so the fraction $p/q$ always makes sense, without a separate clause.
And it puts the sign in the numerator: $-\tfrac34$ is written $\tfrac{-3}{4}$ rather than $\tfrac{3}{-4}$.
Every rational number can be written in many ways ($\tfrac12=\tfrac24=\tfrac{50}{100}$), but this does not
matter: $\Q$ is the set of numbers that can be written in this form at least once.

\begin{proposition}[The number systems are nested][prop:nested]
\[
\N\subseteq\Z\subseteq\Q\subseteq\R. \tag*{\srcin{L1, slide 16}}
\]
\end{proposition}
Every natural number is an integer; every integer $p$ is the rational number $p/1$; and every rational
number is a point on the number line. Equality and the inequalities $\ge,\le,>,<$ are defined on all four
sets as usual, and they agree with one another: $2<3$ means the same whether $2$ and $3$ are thought of as
natural numbers or as real numbers.

\subsection{Why the real numbers are needed}
The slides note that the rational numbers are adequate for \emph{arithmetic}: adding, subtracting,
multiplying or dividing (by a non-zero number) two rational numbers always gives a rational number. The real
numbers are needed for \emph{algebra} (for example, solving $x^2=2$), \emph{geometry} (for example, the ratio of
the circumference of a circle to its diameter, $\pi$) and \emph{limits}. \src{L1, slide 16}

\begin{external}[$\sqrt2$ is not rational]
The equation $x^2=2$ has no solution in $\Q$. Suppose it did: $x=p/q$ with $p\in\Z$, $q\in\N$, and the fraction
in lowest terms (so $p$ and $q$ are not both even). Then $p^2=2q^2$, so $p^2$ is even, and therefore $p$ is even
(the square of an odd number is odd). Write $p=2k$; then $4k^2=2q^2$, so $q^2=2k^2$ is even and $q$ is even.
So $p$ and $q$ are both even, contradicting lowest terms. Hence no rational number squares to $2$. On the
number line there is nevertheless a point at distance $\sqrt2$ from $0$ (the diagonal of a unit square), so the
rational numbers have ``gaps'' that the real numbers fill. \Cref{sec:lub} describes the property of $\R$ that
captures the absence of gaps.
\end{external}

\section{The properties of the real numbers}\label{sec:axioms}

The slides list thirteen properties of the real numbers in three groups: two for the order relation, five
for addition and six for multiplication. In this module they are the axioms: we do not prove them, and every
other fact about real numbers is derived from them.

\subsection{The order relation}

\begin{definition}[Order: O1--O2][def:order]
\begin{description}[style=sameline, leftmargin=3.2em, font=\sffamily\bfseries\color{defcol}]
  \item[O1] For any real numbers $a,b$, exactly one of the following statements is true: either $a>b$, or
  $a=b$, or $a<b$.
  \begin{itemize}
    \item $a<b$ is the same as $b>a$.
    \item We write $a\ge b$ if $a>b$ or $a=b$. We write $a\le b$ if $a<b$ or $a=b$.
  \end{itemize}
  \item[O2] \emph{Transitivity:} if $a>b$ and $b>c$, then $a>c$.
\end{description}
\src{L1, slide 17}
\end{definition}

O1 is called the \term{trichotomy} property (``division into three''). It has two halves: \emph{at least}
one of the three possibilities holds, so any two real numbers can be compared; and \emph{at most} one holds,
so the possibilities exclude each other. The first half is what makes case analyses complete: when a factor
in a product is not zero, it must be positive or negative, and there is no fourth option
(\cref{ex:product-ineq}). The second half is what makes negation simple: ``not $(a>b)$'' is the same as
$a\le b$.

Note that $\ge$ and $\le$ are \emph{defined} from $>$ and $=$; they are not separate primitive notions.

\subsection{Addition}

Addition assigns to any numbers $a,b$ a unique number $a+b$.

\begin{definition}[Addition: A1--A5][def:addition]
For all real numbers $a$, $b$, $c$:
\begin{description}[style=sameline, leftmargin=3.2em, font=\sffamily\bfseries\color{defcol}]
  \item[A1] \emph{Commutativity:} $a+b=b+a$.
  \item[A2] \emph{Associativity:} $a+(b+c)=(a+b)+c$.
  \item[A3] \emph{Property of the number zero:} $a+0=a$.
  \item[A4] \emph{Existence of the additive inverse:} $a+(-a)=0$.
  \item[A5] \emph{Additivity of the order:} $a>b\ \Rightarrow\ a+c>b+c$.
\end{description}
\src{L1, slide 17}
\end{definition}

A1 and A2 say that, in a sum, neither the order of the terms nor the way they are grouped matters, which is
why we may write $a+b+c$ without brackets. A3 says that $0$ changes nothing when added. A4 says that every
number $a$ has a partner $-a$ that cancels it; subtraction is \emph{defined} by $a-b=a+(-b)$. A5 connects
addition to order: adding the same number to both sides of a strict inequality preserves it.

\begin{warning}[A5 is an implication, not an equivalence]
As stated, A5 says only $a>b\Rightarrow a+c>b+c$. The reverse direction is true but needs an argument: from
$a+c>b+c$, apply A5 again with $-c$ to get $(a+c)+(-c)>(b+c)+(-c)$, and simplify with A2--A4 to $a>b$. So
\[
a>b\quad\Longleftrightarrow\quad a+c>b+c \qquad\text{for every real } c.
\]
This is a small theorem, not the axiom itself. The slides use exactly this two-step argument in Example 1 of
\cref{ch:ineq}.
\end{warning}

\subsection{Multiplication}

Multiplication assigns to any numbers $a,b$ a unique number $a\cdot b$ (also written $ab$).

\begin{definition}[Multiplication: M1--M6][def:multiplication]
For all real numbers $a$, $b$, $c$:
\begin{description}[style=sameline, leftmargin=3.2em, font=\sffamily\bfseries\color{defcol}]
  \item[M1] \emph{Commutativity:} $a\cdot b=b\cdot a$.
  \item[M2] \emph{Associativity:} $a\cdot(b\cdot c)=(a\cdot b)\cdot c$.
  \item[M3] \emph{Properties of the number one:} $1\ne0$, and $a\cdot1=a$.
  \item[M4] \emph{Existence of the inverse of a non-zero number:} $a\cdot(1/a)=1$ for $a\ne0$.
  \item[M5] \emph{Distributive law:} $a\cdot(b+c)=a\cdot b+a\cdot c$.
  \item[M6] \emph{Multiplicative properties of the order relation:}
  \begin{itemize}
    \item if $a>b$ and $c>0$, then $a\cdot c>b\cdot c$;
    \item if $a>b$ and $c<0$, then $a\cdot c<b\cdot c$.
  \end{itemize}
\end{description}
\src{L1, slide 18}
\end{definition}

M1--M4 mirror A1--A4. Division is defined by $a/b=a\cdot(1/b)$ for $b\ne0$; there is no $1/0$, which is why
M4 carries the condition $a\ne0$. M5 is the only property that links addition and multiplication: it is what
allows brackets to be multiplied out and common factors to be taken out.

\begin{core}[M6: the sign of the multiplier decides the direction]
Multiplying both sides of an inequality by a \emph{positive} number keeps its direction; multiplying by a
\emph{negative} number reverses it. Multiplying by $0$ destroys it (both sides become $0$). Consequently:
\begin{itemize}
  \item before multiplying or dividing an inequality by anything, establish its sign;
  \item if the sign is unknown (a variable or a parameter), split into the cases positive, zero and negative.
\end{itemize}
This one habit accounts for a large share of the marks in inequality and parameter questions on past ECN115
papers.
\end{core}

\begin{note}[The rational numbers satisfy all thirteen properties]
Every one of O1--O2, A1--A5 and M1--M6 is also true in $\Q$: sums, products, negatives and reciprocals of
rational numbers are rational, and the order behaves in the same way. \src{L1, slide 19} So these properties
cannot, on their own, tell $\R$ apart from $\Q$. Whatever distinguishes the real numbers --- the existence of
$\sqrt2$, the good behaviour of limits --- must come from an additional property (\cref{sec:lub}).
\end{note}

\section{Consequences of the properties}\label{sec:consequences}

Many familiar rules of algebra are not on the list. They are theorems, derived from the list. The slides use
some of them without proof (the study guide points out that Example 2 of \cref{ch:ineq} relies on $0\cdot a=0$);
this section proves the ones needed in the book. The proofs show how theorems are built from axioms, and you
may use the results freely afterwards.

\begin{proposition}[Multiplying by zero][prop:zero-times]
For every real number $a$, $0\cdot a=0$.
\end{proposition}
\emph{Proof.} By A3, $0+0=0$. Multiply by $a$ and use M1 and M5:
$0\cdot a=(0+0)\cdot a=0\cdot a+0\cdot a$. Now add $-(0\cdot a)$ to both sides and use A2, A4 and A3:
$0=0\cdot a$. \hfill$\square$

\begin{proposition}[Zero products][prop:zero-product]
For real numbers $a,b$: if $ab=0$ then $a=0$ or $b=0$. Equivalently, if $a\ne0$ and $b\ne0$ then $ab\ne0$.
\end{proposition}
\emph{Proof.} Suppose $ab=0$ and $a\ne0$. By M4, $1/a$ exists. Then, using M2, M4, M3 and
\cref{prop:zero-times},
$b=1\cdot b=\big((1/a)\cdot a\big)\cdot b=(1/a)\cdot(ab)=(1/a)\cdot0=0$. So if $a\ne0$ then $b=0$; hence
$a=0$ or $b=0$. \hfill$\square$

\begin{proposition}[Signs][prop:signs]
For real numbers $a,b$:
\begin{enumerate}[label=(\roman*)]
  \item $(-1)\cdot a=-a$, and $-(-a)=a$;
  \item $a>0\Leftrightarrow -a<0$;
  \item $a>b\Leftrightarrow a-b>0$;
  \item if $a\ne0$ then $a^2>0$; hence $a^2\ge0$ for every $a$, and in particular $1=1^2>0$;
  \item if $a>0$ then $1/a>0$; if $a<0$ then $1/a<0$;
  \item the product of two numbers with the same sign is positive; the product of two numbers with opposite
  signs is negative.
\end{enumerate}
\end{proposition}
\emph{Proof.} (i) $a+(-1)a=1\cdot a+(-1)a=(1+(-1))a=0\cdot a=0$ by M3, M5, A4 and \cref{prop:zero-times}; so
$(-1)a$ is an additive inverse of $a$. Additive inverses are unique (if $a+b=0$ and $a+c=0$ then
$b=b+(a+c)=(b+a)+c=0+c=c$), so $(-1)a=-a$. Also $-a+a=0$ shows that $a$ is the inverse of $-a$, i.e.\ $-(-a)=a$.

(ii) If $a>0$, add $-a$ to both sides (A5): $0>-a$. Conversely, if $-a<0$, add $a$: $0<a$.

(iii) Add $-b$ to both sides of $a>b$ (A5, A4): $a-b>0$. Conversely add $b$.

(iv) If $a>0$, M6 with $c=a>0$ applied to $a>0$ gives $a\cdot a>0\cdot a=0$. If $a<0$, apply M6 to the
inequality $0>a$ with $c=a<0$: the direction reverses, giving $0\cdot a<a\cdot a$, that is, $a\cdot a>0$.
Either way $a^2>0$.

(v) Let $a>0$. By M4, $a\cdot(1/a)=1>0$. If $1/a$ were $0$ the product would be $0$; if $1/a$ were negative,
M6 applied to $a>0$ with $c=1/a<0$ would give $a\cdot(1/a)<0$. By O1, $1/a>0$. The case $a<0$ is similar.

(vi) If $a>0$ and $b>0$, M6 gives $ab>0\cdot b=0$. If $a<0$ and $b<0$, M6 with $c=b<0$ applied to $0>a$ gives
$0\cdot b<ab$, i.e.\ $ab>0$. If $a>0$ and $b<0$, M6 with $c=b<0$ applied to $a>0$ gives $ab<0$. \hfill$\square$

\begin{note}
You will not be asked to reproduce these derivations in ECN115: the slides say explicitly that, when solving
problems, ``it suffices just to do correct calculations''. They are included so that you can see that every
rule you use rests on the thirteen properties, and so that the reasoning behind sign analysis
(\cref{sec:sign-analysis}) is complete.
\end{note}

\section{The Least Upper Bound property}\label{sec:lub}

\begin{sourcenote}[Non-examinable]
The slides present this property as ``extra material, non-examinable''. It is included because it explains
\emph{why} the real numbers, rather than the rational numbers, are the right setting for limits.
\end{sourcenote}

\begin{definition}[Upper bound and least upper bound][def:lub]
Let $A\subseteq\R$ be non-empty. A number $y\in\R$ is an \term{upper bound} for $A$ if $x\le y$ for all
$x\in A$. A number $z\in\R$ is a \term{least upper bound} for $A$ if
\begin{enumerate}[label=(\roman*)]
  \item $z$ is an upper bound: $x\le z$ for all $x\in A$; and
  \item no smaller number is an upper bound: for every $w<z$ there is $x\in A$ such that $x>w$.
\end{enumerate}
\src{L1, slide 19}
\end{definition}

\begin{principle}[Least Upper Bound property][pr:lub]
If a non-empty subset $A\subseteq\R$ has an upper bound, then $A$ has a least upper bound. \src{L1, slide 19}
\end{principle}

The rational numbers do \emph{not} have this property. The set $A=\{q\in\Q\mid q>0,\ q^2<2\}$ is non-empty
($1\in A$) and bounded above (by $2$), but there is no smallest \emph{rational} upper bound: the least upper
bound would have to be $\sqrt2$, which is not rational. In $\R$ the number $\sqrt2$ exists, and it is the least
upper bound. The Least Upper Bound property is precisely the statement that the real number line has no
gaps; it is what guarantees, in Topic~3, that limits exist when we expect them to.

\section{Intervals}\label{sec:intervals}

Subsets of $\R$ consisting of all numbers between two endpoints occur so often that they have their own
notation.

\begin{definition}[Intervals][def:intervals]
For real numbers $a\le b$:
\begin{align*}
\text{closed interval:}\quad &[a,b]=\{x\in\R\mid a\le x\le b\},\\
\text{open interval:}\quad &(a,b)=\{x\in\R\mid a<x<b\},\\
\text{half-open intervals:}\quad &(a,b]=\{x\in\R\mid a<x\le b\},\qquad [a,b)=\{x\in\R\mid a\le x<b\},\\
\text{infinite intervals:}\quad &(a,\infty)=\{x\in\R\mid a<x\},\qquad (-\infty,b)=\{x\in\R\mid x<b\},\\
\text{half-infinite intervals:}\quad &[a,\infty)=\{x\in\R\mid a\le x\},\qquad (-\infty,b]=\{x\in\R\mid x\le b\},\\
&(-\infty,\infty)=\R.
\end{align*}
\src{L1, slide 20}
\end{definition}

A square bracket means the endpoint is \emph{included} ($\le$); a round bracket means it is \emph{excluded}
($<$). The symbols $\infty$ and $-\infty$ are not real numbers; they only indicate that the interval has no
upper or lower end. They can therefore never be included, and always take a round bracket.

\begin{figure}[ht]
\centering
\begin{tikzpicture}[x=9mm]
  \tikzset{open/.style={circle, draw=accent, fill=white, thick, inner sep=1.6pt},
           closed/.style={circle, draw=accent, fill=accent, thick, inner sep=1.6pt}}
  \foreach \y/\lab/\l/\r in {0/{[1,4]}/closed/closed, -0.9/{(1,4)}/open/open, -1.8/{(1,4]}/open/closed, -2.7/{[1,4)}/closed/open}{
    \draw[-{Stealth}] (-0.5,\y) -- (6.3,\y);
    \draw[accent, line width=2.2pt] (1,\y) -- (4,\y);
    \node[\l] at (1,\y) {}; \node[\r] at (4,\y) {};
    \node[anchor=west, font=\small] at (6.6,\y) {$\lab$};
  }
  \draw[-{Stealth}] (-0.5,-3.6) -- (6.3,-3.6);
  \draw[accent, line width=2.2pt, -{Stealth}] (1,-3.6) -- (6.2,-3.6);
  \node[open] at (1,-3.6) {};
  \node[anchor=west, font=\small] at (6.6,-3.6) {$(1,\infty)$};
  \draw[-{Stealth}] (-0.5,-4.5) -- (6.3,-4.5);
  \draw[accent, line width=2.2pt, {Stealth}-] (-0.4,-4.5) -- (4,-4.5);
  \node[closed] at (4,-4.5) {};
  \node[anchor=west, font=\small] at (6.6,-4.5) {$(-\infty,4]$};
  \foreach \x in {0,...,6} \draw (\x,-4.9) -- (\x,-5.05) node[below, font=\scriptsize] {\x};
\end{tikzpicture}
\caption{Intervals on the number line. A filled dot marks an included endpoint (square bracket), an open dot an
excluded endpoint (round bracket). An arrow shows that the interval continues without end.}
\label{fig:intervals}
\end{figure}

\begin{example}[Writing sets as intervals]
Write each set as an interval or a union of intervals: (a) $\{x\in\R\mid -2\le x<7\}$;
(b) $\{x\in\R\mid x>3\}$; (c) $\{x\in\R\mid x\le 0\text{ or }x>5\}$; (d) $\{x\in\R\mid x\ne 1\}$.
\solution
(a) $[-2,7)$. (b) $(3,\infty)$. (c) $(-\infty,0]\cup(5,\infty)$. (d) $(-\infty,1)\cup(1,\infty)$, which can also be
written $\R\setminus\{1\}$.
\end{example}

\begin{external}[Why the type of endpoint matters]
$[0,1)$ has a smallest element ($0$) but no largest: for any $x<1$ in the interval, $(x+1)/2$ is a larger element
of the interval. $(0,1]$ has a largest element ($1$) but no smallest. In optimisation (Topic~4) this decides
whether a function defined on the interval can attain a maximum or a minimum, which is why a question may be
posed on a half-open interval such as $(0,c]$.
\end{external}

% --------------------------------------------------------------------
\section{Chapter review}
% --------------------------------------------------------------------
\subsection*{Summary}
$\N=\{1,2,3,\dots\}$, with the successor property; $\Z$ adds zero and the negatives; $\Q$ consists of fractions
$p/q$ with $p\in\Z$, $q\in\N$; and $\R$ is, informally, the number line. They are nested:
$\N\subseteq\Z\subseteq\Q\subseteq\R$. The rationals suffice for arithmetic, but solving $x^2=2$, measuring a
circle and taking limits need $\R$. The real numbers satisfy the order properties O1 (trichotomy) and O2
(transitivity), the addition properties A1--A5 and the multiplication properties M1--M6; in ECN115 these are
the axioms. M6 --- a positive multiplier keeps an inequality, a negative one reverses it --- is the most
important for calculation. $\Q$ satisfies all thirteen properties too; the additional Least Upper Bound property
(non-examinable) is what distinguishes $\R$. Intervals name the subsets of $\R$ between two endpoints.

\subsection*{Key definitions and properties}
\begin{itemize}
  \item $\N=\{1,2,3,\dots\}$; \ $\Z=\{\dots,-1,0,1,\dots\}$; \ $\Q=\{p/q: p\in\Z,\ q\in\N\}$; \ $\R$: the number line.
  \item O1: exactly one of $a>b$, $a=b$, $a<b$. \ O2: $a>b$, $b>c\Rightarrow a>c$.
  \item A1--A4: $a+b=b+a$; $a+(b+c)=(a+b)+c$; $a+0=a$; $a+(-a)=0$. \ A5: $a>b\Rightarrow a+c>b+c$.
  \item M1--M5: $ab=ba$; $a(bc)=(ab)c$; $1\ne0$, $a\cdot1=a$; $a\cdot(1/a)=1$ $(a\ne0)$; $a(b+c)=ab+ac$.
  \item M6: $a>b$, $c>0\Rightarrow ac>bc$; \ $a>b$, $c<0\Rightarrow ac<bc$.
  \item Intervals: $[a,b]$, $(a,b)$, $(a,b]$, $[a,b)$, $(a,\infty)$, $[a,\infty)$, $(-\infty,b)$, $(-\infty,b]$, $\R$.
\end{itemize}

\subsection*{Connections}
The axioms justify every step of the inequality solving in \cref{ch:ineq}; the sign rules of
\cref{prop:signs} drive the sign analysis there. The successor property underlies induction (\cref{ch:ind}).
Intervals are the natural form of answers to inequalities and the domains of functions (\cref{ch:func}). The
Least Upper Bound property is the reason limits (\cref{ch:limits}) are studied in $\R$ rather than $\Q$.

\subsection*{Common mistakes}
\begin{itemize}
  \item Including $0$ in $\N$.
  \item Forgetting to reverse an inequality after multiplying or dividing by a negative number.
  \item Dividing by a quantity that might be zero.
  \item Putting a square bracket next to $\infty$, or the wrong bracket for a strict inequality.
  \item Assuming A5 gives an equivalence without the second step.
\end{itemize}

\subsection*{Exam checklist}
\begin{checklist}
  \item Write down $\N$, $\Z$, $\Q$, $\R$ and the inclusions between them; state the successor property.
  \item Explain why $\R$ is needed in addition to $\Q$.
  \item State O1, O2, A1--A5 and M1--M6, and say which property justifies a given step.
  \item State both clauses of M6 exactly.
  \item Write any interval in set-builder notation and vice versa.
\end{checklist}

% --------------------------------------------------------------------
\section{Practice questions}
% --------------------------------------------------------------------
\pq[basic] Decide whether each number belongs to $\N$, $\Z$, $\Q$, $\R$ (list all that apply):
$5$, $-3$, $0$, $\tfrac{-7}{2}$, $0.25$, $\sqrt9$, $\pi$.

\pq[basic] Write in set-builder notation: (a) $(-3,5]$; (b) $[2,\infty)$; (c) $(-\infty,0)\cup[1,2]$.

\pq[conceptual] Explain why the definition of $\Q$ uses $q\in\N$ rather than $q\in\Z$.

\pq[mathematical] Starting from $a>b$, and naming the property used at each step, show that $-a<-b$.

\pq[mathematical] Using the properties of the real numbers, prove that if $a>b$ and $c>d$ then $a+c>b+d$.

\pq[conceptual] Which of the thirteen properties O1--M6 fail for the natural numbers $\N$ (with the usual
addition and multiplication)? Give a reason for each.

\pq[multi-step] Let $a,b$ be real numbers with $0<a<b$. Prove that $a^2<b^2$ and that $1/a>1/b$.

\pq[exam-style] (a) State both clauses of the multiplicative property of the order relation (M6).
(b) A student argues: ``$x>y$, so multiplying both sides by $x$ gives $x^2>xy$.'' Give conditions on $x$
under which the conclusion is guaranteed, and give an example with $x>y$ for which $x^2>xy$ is false.

% --------------------------------------------------------------------
\section{Solutions to practice questions}
% --------------------------------------------------------------------
\pqsol{3.1}
$5$: $\N,\Z,\Q,\R$. \ $-3$: $\Z,\Q,\R$. \ $0$: $\Z,\Q,\R$ (not $\N$). \ $\tfrac{-7}{2}$: $\Q,\R$. \
$0.25=\tfrac14$: $\Q,\R$. \ $\sqrt9=3$: $\N,\Z,\Q,\R$. \ $\pi$: $\R$ only.

\pqsol{3.2}
(a) $\{x\in\R\mid -3<x\le5\}$. (b) $\{x\in\R\mid x\ge2\}$. (c) $\{x\in\R\mid x<0\text{ or }1\le x\le2\}$.

\pqsol{3.3}
Since $0\notin\N$, the condition $q\in\N$ excludes a zero denominator automatically, so $p/q$ is always
defined. It also fixes the sign in the numerator. With $q\in\Z$ a separate condition $q\ne0$ would be needed.

\pqsol{3.4}
Add $(-a)+(-b)$ to both sides of $a>b$ (A5): $a+((-a)+(-b))>b+((-a)+(-b))$. Rearranging with A1 and A2,
the left side is $(a+(-a))+(-b)=0+(-b)=-b$ (A4, A1, A3) and the right side is $(b+(-b))+(-a)=-a$. Hence
$-b>-a$, i.e.\ $-a<-b$. (Alternatively: multiply by $-1<0$ using M6 and \cref{prop:signs}(i).)

\pqsol{3.5}
From $a>b$, A5 with $c$ gives $a+c>b+c$. From $c>d$, A5 with $b$ gives $c+b>d+b$, i.e.\ $b+c>b+d$ (A1).
By transitivity (O2), $a+c>b+d$.

\pqsol{3.6}
A3 fails: there is no $0$ in $\N$, so ``$a+0=a$'' has no meaning inside $\N$. A4 fails: $-1\notin\N$, so $1$
has no additive inverse in $\N$. M4 fails: $2\in\N$ but $1/2\notin\N$. The order properties O1, O2, A5 and M6
(with $c\in\N$; the second clause never applies since no natural number is negative), and A1, A2, M1--M3, M5
continue to hold.

\pqsol{3.7}
Since $b>a$ and $a>0$, M6 gives $b\cdot a>a\cdot a$; since $b>a$ and $b>0$, M6 gives $b\cdot b>a\cdot b$. By M1 and
O2, $b^2>ab>a^2$. For the reciprocals: $a>0$ and $b>0$ give $ab>0$, so $1/(ab)>0$ by \cref{prop:signs}(v).
Multiply $b>a$ by $1/(ab)>0$ (M6): $b/(ab)>a/(ab)$, i.e.\ $1/a>1/b$.

\pqsol{3.8}
(a) If $a>b$ and $c>0$, then $a\cdot c>b\cdot c$. If $a>b$ and $c<0$, then $a\cdot c<b\cdot c$.
(b) The step is M6 with multiplier $x$, so it is guaranteed when $x>0$. If $x<0$ the direction reverses and
$x^2<xy$; if $x=0$ both sides are $0$. Example: $x=-1$, $y=-2$ gives $x>y$, but $x^2=1$ and $xy=2$, so
$x^2>xy$ is false.
